\chapter{Situação Atual}
\label{cap:situacao-atual}

Este capítulo apresenta o contexto necessário para compreender o problema tratado neste trabalho. A Seção~\ref{sec:conceitos} define os conceitos e a terminologia básica. As Seções~\ref{sec:previsao} e~\ref{sec:alocacao} descrevem, respectivamente, as abordagens existentes para previsão de carga e para alocação de recursos em nuvem, destacando suas limitações. A Seção~\ref{sec:lacuna} compara os trabalhos mais próximos desta proposta, e a Seção~\ref{sec:tecnologias} apresenta as tecnologias utilizadas.

\section{Conceitos e terminologia}
\label{sec:conceitos}

\subsection{Elasticidade, carga de trabalho e séries temporais}
\label{subsec:carga}

Uma das características essenciais da computação em nuvem é a elasticidade rápida, isto é, a capacidade de provisionar e liberar recursos rapidamente, em alguns casos de forma automática, acompanhando a demanda \cite{mell2011nist}. \citeonline{herbst2013elasticity} definem elasticidade como o grau em que um sistema se adapta a mudanças na carga de trabalho, provisionando e liberando recursos de forma autônoma, de modo que, a cada instante, os recursos disponíveis correspondam o mais fielmente possível à demanda corrente. Essa definição deixa claro que a elasticidade tem duas falhas possíveis: oferecer recursos de mais ou recursos de menos.

A carga de trabalho (\textit{workload}) de um \textit{cluster} é o conjunto de demandas submetidas a ele, como tarefas em lote, serviços de longa duração e requisições de usuários. Do ponto de vista da infraestrutura, essa carga se manifesta como utilização de recursos, principalmente processador e memória. Quando essa utilização é agregada em intervalos regulares de duração $\Delta$, obtém-se uma série temporal $y_1, y_2, \dots, y_t$, em que $y_t$ é a utilização média (ou máxima) no intervalo $t$. Séries de carga em nuvem costumam combinar tendência, padrões periódicos (por exemplo, diários e semanais), ruído e picos abruptos, o que torna sua previsão difícil \cite{xu2022esdnn}.

Prever a carga significa estimar $\hat{y}_{t+h}$, o valor da série $h$ intervalos à frente, a partir de uma janela com as $w$ observações mais recentes. O parâmetro $h$ é o horizonte de previsão e $w$ é o tamanho da janela de entrada. Quando é preciso distinguir o intervalo em que a previsão é feita do intervalo previsto, o primeiro é chamado de origem da previsão e denotado por $\tau$: uma previsão feita na origem $\tau$ só pode usar observações até $\tau$. A previsão é univariada quando usa apenas a própria série e multivariada quando incorpora outras séries correlacionadas, como prever o uso de processador a partir do histórico de processador e de memória \cite{xu2022esdnn}.

Pesquisas sobre o tema dependem de rastros (\textit{traces}) públicos de \textit{clusters} reais, que registram eventos e uso de recursos de máquinas, tarefas e contêineres. O rastro mais recente do Google cobre oito células do gerenciador Borg durante o mês de maio de 2019 \cite{wilkes2020google}. O rastro de 2018 da Alibaba cobre cerca de 4.000 máquinas de um \textit{cluster} de produção ao longo de oito dias, com serviços \textit{online} e tarefas em lote compartilhando as mesmas máquinas, e inclui tabelas de uso de recursos por máquina e por contêiner \cite{alibaba2018trace}. Como esses rastros registram eventos e medições em granularidade fina, transformá-los em séries temporais regulares é uma etapa de pré-processamento necessária, que faz parte deste trabalho.

A granularidade dessa agregação varia entre os trabalhos, mas intervalos de cinco minutos são os mais usados com esses rastros. \citeonline{rossi2023forecasting} agregam os rastros do Google e da Alibaba em intervalos de cinco minutos e preveem a demanda dez minutos à frente, horizonte que consideram suficiente para a maioria das aplicações de alocação. O esDNN também adota cinco minutos como intervalo de previsão e de decisão na simulação de escalonamento \cite{xu2022esdnn}. Trabalhos com outros tipos de dado usam escalas diferentes, como uma hora para requisições a servidores \textit{web} \cite{calheiros2015arima} ou um minuto para dados de produção de um provedor \cite{pan2023magicscaler}. O ponto comum é que o intervalo deve ser longo o bastante para que uma nova instância possa ser iniciada dentro dele \cite{calheiros2015arima}.

\subsection{Acordos de nível de serviço}
\label{subsec:sla}

A qualidade de um serviço em nuvem é expressa por três conceitos relacionados \cite{beyer2016sre}. Um indicador de nível de serviço (\textit{Service Level Indicator}, SLI) é uma medida quantitativa de algum aspecto do serviço, como latência, disponibilidade ou fração de requisições atendidas com sucesso. Um objetivo de nível de serviço (\textit{Service Level Objective}, SLO) é um valor-alvo ou uma faixa aceitável para um SLI, por exemplo, atender 99\% das requisições dentro de um limite de tempo. Por fim, um acordo de nível de serviço (\textit{Service Level Agreement}, SLA) é o contrato entre provedor e cliente que formaliza esses objetivos e define as consequências, geralmente financeiras, de não cumpri-los.

Na literatura de alocação de recursos, a violação de SLA é frequentemente aproximada por medidas derivadas da utilização, já que os rastros públicos não registram a latência percebida pelos usuários. \citeonline{beloglazov2012optimal}, por exemplo, medem a fração do tempo em que os servidores ativos operam com utilização máxima, situação em que o desempenho das aplicações é presumivelmente degradado. \citeonline{rossi2023forecasting} consideram atendida a demanda que fica abaixo da capacidade prevista.

Este trabalho adota uma aproximação do mesmo tipo. O SLI é medido por intervalo: o intervalo $t$ é atendido integralmente quando a capacidade disponível cobre a demanda real, isto é, quando $d_t \leq k \cdot n_t$, na notação da Equação~\ref{eq:problema}, e há violação no caso contrário. O SLO correspondente é que no máximo uma fração $\varepsilon$ dos intervalos fique sem atendimento integral, que é exatamente a restrição daquela equação. A fração da demanda que deixa de ser atendida é acompanhada como métrica complementar (Seção~\ref{subsec:metricas}). Essa simplificação não captura efeitos como filas ou degradação gradual de desempenho perto da saturação, mas permite comparar estratégias de alocação de forma consistente e reproduzível.

\subsection{Provisionamento reativo e preditivo}
\label{subsec:provisionamento}

O escalonamento automático (\textit{autoscaling}) é o mecanismo que ajusta os recursos alocados a uma aplicação sem intervenção humana. Ele pode ser horizontal, quando altera o número de instâncias (máquinas virtuais ou contêineres), ou vertical, quando altera a capacidade de cada instância. Este trabalho trata apenas do escalonamento horizontal.

\citeonline{lorido2014autoscaling} classificam as técnicas de escalonamento automático em cinco famílias: regras baseadas em limiares, aprendizado por reforço, teoria de filas, teoria de controle e análise de séries temporais. Do ponto de vista deste trabalho, a distinção mais importante é entre abordagens reativas e preditivas. No provisionamento reativo, a decisão no instante $t$ é tomada a partir do que já foi observado, normalmente a utilização corrente. No provisionamento preditivo, a decisão é tomada a partir de uma estimativa $\hat{d}_{t+h}$ da demanda futura. Na alocação, a série prevista é a própria demanda, isto é, $y_t = d_t$.

A desvantagem da abordagem reativa é o atraso. Entre a detecção de um aumento de carga e a disponibilidade de uma nova instância há o atraso de inicialização, que não é desprezível e varia entre provedores e configurações \cite{mao2012startup}. Durante esse período, a aplicação opera com capacidade insuficiente. Por isso, a decisão de alocação precisa ser tomada com antecedência. Neste trabalho, a decisão tomada com as observações disponíveis até o intervalo $\tau$ passa a valer no intervalo $\tau + L$, em que $L$ é a antecedência da decisão (\textit{lead time}), igual a um intervalo mais o atraso de inicialização, também medido em intervalos. Sem atraso de inicialização, $L = 1$ e a decisão vale já no intervalo seguinte. As duas abordagens estão sujeitas à mesma antecedência: a reativa decide em $\tau$ com a utilização observada em $\tau$, e a preditiva tenta antecipar a demanda do intervalo $\tau + L$, de modo que a instância já esteja disponível quando a carga chegar. Provedores comerciais já oferecem modos de escalonamento preditivo, como o do Amazon EC2 Auto Scaling \cite{aws2024predictive}, o que reforça a relevância prática do problema.

Qualquer que seja a abordagem, cada decisão pode errar em duas direções. Há superprovisionamento quando a capacidade alocada excede a demanda ($k \cdot n_t > d_t$): o serviço é atendido, mas a capacidade ociosa $k \cdot n_t - d_t$ representa custo sem retorno. Há subprovisionamento quando a capacidade é menor que a demanda ($k \cdot n_t < d_t$): o custo é menor, mas parte da demanda não é atendida e o SLA pode ser violado. O objetivo da alocação é equilibrar essas duas falhas, e a qualidade da previsão determina o quanto esse equilíbrio pode ser alcançado.

\subsection{Métricas de avaliação}
\label{subsec:metricas}

As métricas usadas na literatura se dividem em dois grupos, e a distinção entre eles é central para este trabalho.

\textbf{Métricas de previsão.} Medem o quanto a previsão se aproxima do valor real, independentemente de como ela será usada. Para $N$ previsões $\hat{y}_i$ e valores reais $y_i$, as mais comuns são o erro quadrático médio (MSE), sua raiz (RMSE), o erro absoluto médio (MAE) e o erro percentual absoluto médio (MAPE):

\begin{align}
\mathrm{MSE}  &= \frac{1}{N}\sum_{i=1}^{N}(y_i - \hat{y}_i)^2, &
\mathrm{RMSE} &= \sqrt{\mathrm{MSE}}, \label{eq:mse-rmse}\\
\mathrm{MAE}  &= \frac{1}{N}\sum_{i=1}^{N}\lvert y_i - \hat{y}_i \rvert, &
\mathrm{MAPE} &= \frac{100\%}{N}\sum_{i=1}^{N}\left\lvert \frac{y_i - \hat{y}_i}{y_i} \right\rvert. \label{eq:mae-mape}
\end{align}

O MAPE é indefinido quando algum $y_i$ é zero e instável quando $y_i$ é próximo de zero, situação comum em períodos de baixa utilização. Neste trabalho, ele é calculado apenas sobre os pontos em que $y_i$ não é próximo de zero, e a fração de pontos mantida, chamada de cobertura da máscara, é informada junto com o valor. Além disso, as quatro métricas são simétricas: superestimar e subestimar a demanda pela mesma quantidade produz o mesmo erro. Para a alocação, porém, esses erros têm consequências diferentes, já que a superestimação gera custo e a subestimação gera violações. Por isso, um modelo com menor RMSE não é necessariamente o que leva às melhores decisões de alocação. Para captar essa assimetria, são usadas também duas medidas que distinguem o sentido do erro: o viés, que é o erro médio com sinal, e a taxa de subestimação, que é a fração das previsões abaixo do valor real:

\begin{align}
\text{viés} &= \frac{1}{N}\sum_{i=1}^{N}(\hat{y}_i - y_i), &
\text{taxa de subestimação} &= \frac{1}{N}\sum_{i=1}^{N}\mathbf{1}\left[\, \hat{y}_i < y_i \,\right]. \label{eq:vies}
\end{align}

\noindent Um viés negativo indica que o previsor tende a subestimar a demanda e, portanto, a provocar violações quando usado na alocação sem margem.

\textbf{Métricas de alocação.} Medem o resultado das decisões tomadas a partir da previsão. Para uma simulação com $T$ intervalos, as métricas adotadas neste trabalho são:

\begin{align}
C &= \sum_{t=1}^{T} c \cdot n_t,
  && \text{(custo operacional total)} \label{eq:custo}\\
V &= \frac{1}{T}\sum_{t=1}^{T} \mathbf{1}\left[\, d_t > k \cdot n_t \,\right],
  && \text{(fração de intervalos com violação)} \label{eq:violacao}\\
U &= \frac{\sum_{t} \max(0,\, d_t - k \cdot n_t)}{\sum_{t} d_t},
  && \text{(fração da demanda não atendida)} \label{eq:nao-atendida}\\
O &= \frac{\sum_{t} \max(0,\, k \cdot n_t - d_t)}{\sum_{t} k \cdot n_t},
  && \text{(fração da capacidade ociosa)} \label{eq:ociosa}\\
A &= \sum_{t=2}^{T} \mathbf{1}\left[\, n_t \neq n_{t-1} \,\right].
  && \text{(número de alterações em $n_t$)} \label{eq:alteracoes}
\end{align}

\noindent O número de alterações $A$ é acompanhado porque políticas que oscilam com frequência são indesejáveis na prática. Essas métricas correspondem diretamente aos termos da formulação da Equação~\ref{eq:problema}: $C$ é a função objetivo e $V$ é o lado esquerdo da restrição.

Para comparar políticas, o custo é normalizado pelo de um oráculo, uma política de referência que conhece a demanda futura e aloca, em cada intervalo, o menor número de instâncias que a atende, $\lceil d_t / k \rceil$:

\begin{equation}
\label{eq:custo-normalizado}
C_{\mathrm{norm}} = \frac{C}{C_{\text{oráculo}}}.
\end{equation}

\noindent Como nenhuma alocação mais barata que a do oráculo atende toda a demanda, $C_{\text{oráculo}}$ é o piso do custo sem violações. Uma política que viola pode custar menos que o oráculo, e por isso $C_{\mathrm{norm}}$ deve ser lido sempre junto com $V$.

\section{Abordagens para previsão de carga}
\label{sec:previsao}

\subsection{Modelos estatísticos}
\label{subsec:estatisticos}
Os métodos mais simples de previsão servem como referência para todos os demais \cite{hyndman2018forecasting}. O método ingênuo, ou de persistência, assume que o próximo valor será igual ao último observado, $\hat{y}_{t+1} = y_t$. A média móvel usa a média das $w$ últimas observações. O sazonal ingênuo repete o valor observado um período sazonal antes, $\hat{y}_{t+h} = y_{t+h-s}$, em que $s$ é o período da sazonalidade, medido em intervalos, e $h \leq s$. A suavização exponencial simples prevê uma média ponderada das observações passadas, com pesos que decaem geometricamente com a idade da observação. O método de Holt-Winters estende essa ideia com componentes de tendência e de sazonalidade, o que o torna adequado a séries com padrões periódicos, como os ciclos diários comuns em cargas de nuvem \cite{hyndman2018forecasting}.

A família ARIMA (\textit{AutoRegressive Integrated Moving Average}), sistematizada pela metodologia de Box e Jenkins \cite{box2015time}, é o modelo estatístico mais usado em previsão de séries temporais. Um modelo ARIMA$(p,d,q)$ aplica $d$ diferenciações para tornar a série estacionária e descreve o resultado como combinação linear de $p$ valores passados (componente autorregressiva) e de $q$ erros de previsão passados (componente de média móvel). Usando o operador de defasagem $B$, definido por $B\,y_t = y_{t-1}$, o modelo é escrito como:

\begin{equation}
\label{eq:arima}
\left(1 - \sum_{i=1}^{p}\phi_i B^i\right)(1 - B)^d\, y_t
= \beta_0 + \left(1 + \sum_{j=1}^{q}\theta_j B^j\right)a_t,
\end{equation}

\noindent em que $\phi_i$ e $\theta_j$ são os coeficientes estimados, $\beta_0$ é uma constante e $a_t$ é um ruído branco. A variante sazonal, SARIMA (\textit{Seasonal ARIMA}), acrescenta termos equivalentes defasados pelo período da sazonalidade. A escolha das ordens $p$, $d$ e $q$ pode ser feita pela análise das funções de autocorrelação ou por critérios de informação, como o critério de informação de Akaike (AIC) \cite{hyndman2018forecasting}.

No contexto de nuvem, \citeonline{calheiros2015arima} aplicaram um modelo ARIMA a rastros reais de requisições à Wikipédia em inglês, agregados por hora, para apoiar o provisionamento proativo de recursos de provedores de \textit{software} como serviço. O modelo prevê a demanda da hora seguinte, e um módulo de desempenho converte essa previsão no número de máquinas virtuais necessárias. Além da precisão da previsão, que chegou a 91\% em média, os autores avaliaram em simulação, com o CloudSim \cite{calheiros2011cloudsim}, o efeito das previsões sobre a utilização dos recursos e a qualidade de serviço (QoS). Um ponto especialmente relevante para este trabalho é que eles repetiram a simulação usando, em vez do valor previsto, os limites inferiores e superiores dos intervalos de confiança de 80\% e 95\%. Usar os limites superiores reduziu as requisições rejeitadas e as violações de QoS, ao custo de mais horas de máquina virtual, o que mostra de forma concreta o equilíbrio entre custo e violações controlado por um ajuste aplicado depois da previsão. A avaliação, contudo, envolve um único modelo de previsão.

\textbf{Limitações.} Os modelos estatísticos tradicionais assumem relações lineares e, no caso do ARIMA, estacionariedade após a diferenciação. Cargas de nuvem frequentemente violam essas hipóteses, com mudanças abruptas de regime e picos que não seguem um padrão linear. Esses modelos também são essencialmente univariados em sua forma básica e precisam ser reajustados quando o comportamento da série muda. Por outro lado, eles exigem poucos dados, são baratos de ajustar e têm parâmetros interpretáveis. Por isso, devem estar presentes em qualquer comparação: uma rede neural só se justifica se superar essas alternativas mais simples sob as mesmas condições.

\subsection{Redes neurais recorrentes}
\label{subsec:recorrentes}

Para aplicar uma rede neural à previsão, a série é transformada em um problema de aprendizado supervisionado por meio de uma janela deslizante: cada exemplo tem como entrada as $w$ observações mais recentes e como saída o valor a ser previsto. Para vários horizontes, o modelo pode ser aplicado de forma recursiva, realimentando as próprias previsões e acumulando erro, ou de forma direta, com uma saída para cada horizonte \cite{taieb2012strategies}.

As redes neurais recorrentes (RNN) leem a janela em ordem e mantêm um estado interno que funciona como memória do que já foi visto, o que as torna adequadas a séries temporais. Em sequências longas, porém, o treinamento tem dificuldade de aprender relações entre observações distantes \cite{bengio1994learning}. A LSTM (\textit{Long Short-Term Memory}) resolve esse problema com portas que aprendem o que guardar e o que esquecer \cite{hochreiter1997lstm}, e a GRU (\textit{Gated Recurrent Unit}) simplifica essa ideia com apenas duas portas \cite{cho2014gru}, sendo mais barata de treinar e com desempenho comparável \cite{chung2014empirical}.

O esDNN \cite{xu2022esdnn} combina uma GRU modificada com camadas convolucionais para prever o uso de processador e de memória nos rastros da Alibaba e do Google. Nos experimentos, a GRU pura é competitiva em horizontes curtos, e o esDNN se destaca em horizontes longos. A avaliação, no entanto, apoia-se em métricas de erro simétricas, que não distinguem superestimação de subestimação (Seção~\ref{subsec:metricas}).

\textbf{Limitações.} Redes recorrentes exigem mais dados e mais ajuste do que os modelos estatísticos, e, como qualquer previsor treinado para minimizar um erro médio, tendem a suavizar os picos de demanda. \citeonline{agullo2025enhancing} mostram que previsores tradicionais e profundos falham justamente nesses picos, sem que as métricas usuais revelem a falha, e que ajustar a saída da previsão melhora a antecipação dos picos ao custo de mais superprovisionamento. Isso reforça que o erro médio não basta para avaliar um previsor usado na alocação.

\subsection{Modelos com quantificação de incerteza}
\label{subsec:incerteza}

\citeonline{rossi2023forecasting} propõem dois modelos profundos com quantificação de incerteza, uma rede neural bayesiana (HBNN) e uma LSTM com saída distribucional (LSTMD), em versões univariada e bivariada, que preveem não apenas um valor, mas uma distribuição para a demanda futura de processador e de memória. Os experimentos usam rastros do Google e da Alibaba e comparam os modelos com uma LSTM determinística. Os autores mostram que as métricas de erro pontual não distinguem bem os modelos, enquanto métricas de nível de serviço, calculadas sobre o limite superior do intervalo de confiança, evidenciam a vantagem de modelar a incerteza: a taxa de sucesso, fração dos intervalos em que a demanda real fica abaixo desse limite, e o total de recursos previstos, soma dos limites em relação à demanda real. O trabalho também investiga transferência de aprendizado e conclui que ela funciona razoavelmente dentro de um mesmo provedor, mas se degrada entre provedores diferentes.

\textbf{Limitações.} O efeito da incerteza é medido por métricas calculadas sobre a própria previsão, sem acoplar os modelos a um mecanismo de alocação, e os próprios autores apontam essa integração como trabalho futuro \cite{rossi2023forecasting}. Além disso, as métricas de nível de serviço supõem que a capacidade possa ser ajustada de forma contínua, enquanto, na prática, os recursos são alocados em instâncias discretas, com atraso de inicialização. Este trabalho dá esse passo ao aplicar a previsão a uma política com instâncias discretas e medir diretamente as métricas de alocação. Ainda assim, o trabalho sugere uma ideia simples e útil: provisionar pelo limite superior de um intervalo de confiança, e não pelo valor central, é uma forma de trocar custo por menos violações. A margem de segurança $m$ apresentada na Seção~\ref{subsec:alocacao-preditiva} é a versão determinística dessa mesma ideia e é a adotada neste trabalho.

%======================================================================
\section{Abordagens para alocação de recursos}
\label{sec:alocacao}

\subsection{Heurísticas e políticas reativas}
\label{subsec:heuristicas}
As políticas baseadas em limiares são a forma mais difundida de escalonamento automático, sobretudo pela simplicidade \cite{lorido2014autoscaling}. Na forma mais comum, define-se um limiar superior $\theta_{\mathrm{sup}}$ e um inferior $\theta_{\mathrm{inf}}$ para a utilização média das instâncias: se a utilização permanece acima de $\theta_{\mathrm{sup}}$ por um certo número de intervalos, uma instância é adicionada; se permanece abaixo de $\theta_{\mathrm{inf}}$, uma instância é removida. Um período de espera (\textit{cooldown}) após cada ação impede que o sistema oscile.

Uma variação proporcional é usada pelo \textit{Horizontal Pod Autoscaler} (HPA) do Kubernetes, que calcula o número desejado de réplicas a partir da razão entre a métrica observada $u_t$ e um valor-alvo $u^{*}$ \cite{kubernetes2024hpa}:

\begin{equation}
\label{eq:hpa}
n_{t+1} = \left\lceil n_t \cdot \frac{u_t}{u^{*}} \right\rceil.
\end{equation}

\noindent O mecanismo inclui ainda uma faixa de tolerância, dentro da qual nenhuma ação é tomada, e uma janela de estabilização para reduções, ambas destinadas a evitar oscilações. Outro exemplo de política reativa é a usada como referência pelo esDNN, que define o número de máquinas ativas como a média do número de máquinas dos cinco intervalos anteriores \cite{xu2022esdnn}.

Um problema relacionado, mas distinto, é o posicionamento das instâncias nos servidores físicos. Ele é tratado como uma variante do problema de empacotamento (\textit{bin packing}), que é NP-difícil, e por isso costuma ser resolvido com heurísticas como \textit{First-Fit} e \textit{Best-Fit Decreasing}. \citeonline{beloglazov2012optimal}, por exemplo, adaptam esta última para consolidar máquinas virtuais e reduzir o consumo de energia. Meta-heurísticas como algoritmos genéticos, otimização por enxame de partículas e colônia de formigas exploram o espaço de soluções de forma mais ampla, ao custo de maior tempo de computação \cite{bodra2025ml}. Como definido na Seção~\ref{sec:escopo}, o posicionamento está fora do escopo deste trabalho, que decide apenas quantas instâncias manter ativas.

\textbf{Limitações.} Por construção, as políticas reativas só respondem depois que a mudança de carga já ocorreu, e o atraso de inicialização das instâncias aumenta essa defasagem. Os limiares precisam ser ajustados manualmente para cada aplicação: limiares agressivos causam oscilações, e limiares conservadores mantêm capacidade ociosa. Ainda assim, essas políticas são simples, previsíveis e amplamente usadas, o que faz de uma política reativa uma referência obrigatória na avaliação deste trabalho. A referência adotada é a variante proporcional do HPA (Equação~\ref{eq:hpa}), com a faixa de tolerância e a janela de estabilização descritas acima.

\subsection{Alocação guiada por previsão}
\label{subsec:alocacao-preditiva}

Na alocação guiada por previsão, a decisão é tomada em duas etapas desacopladas. Primeiro, um modelo estima a demanda $\hat{d}_{t+h}$ para o intervalo em que a decisão terá efeito. Em seguida, uma regra converte essa estimativa em número de instâncias, normalmente dividindo a demanda prevista pela capacidade de uma instância e acrescentando uma margem de segurança $m \geq 0$, com um piso de $n_{\min}$ instâncias:

\begin{equation}
\label{eq:alocacao-margem}
n_{t+h} = \max\left(n_{\min},\, \left\lceil \frac{(1 + m)\, \hat{d}_{t+h}}{k} \right\rceil\right).
\end{equation}

\noindent Na alocação, $t$ é a origem $\tau$ da previsão e o horizonte é a antecedência da decisão, $h = L$, de modo que a previsão se refira ao intervalo em que a decisão vale e cubra o atraso de inicialização das instâncias \cite{mao2012startup}, como discutido na Seção~\ref{subsec:provisionamento}. A margem controla o equilíbrio entre custo e violações: margens maiores reduzem $V$ e aumentam $C$. Escolhê-la, calibrando em um conjunto de validação, é a forma prática de buscar o menor custo que respeita o limite $\varepsilon$ da Equação~\ref{eq:problema}. Quando o modelo fornece uma distribuição, a margem pode ser substituída por um quantil, como fazem \citeonline{rossi2023forecasting} e \citeonline{calheiros2015arima} ao provisionar pelo limite superior de um intervalo de confiança.

O MagicScaler \cite{pan2023magicscaler} leva essa ideia mais longe: combina uma previsão com incerteza a uma otimização estocástica que minimiza o custo sujeito a um limite de demanda não atendida e, em \textit{clusters} da Alibaba Cloud, reduziu substancialmente a perda de qualidade de serviço em relação a uma política reativa, com pequeno aumento de custo. Seus dados, porém, são privados, e o mecanismo de decisão é aplicado apenas ao próprio previsor, sem comparação com outros modelos sob a mesma regra.

A principal vantagem do desacoplamento é a modularidade: com a regra de decisão fixa, é possível trocar apenas o modelo de previsão e atribuir a ele as diferenças de custo e de violações. É essa propriedade que este trabalho explora para comparar redes neurais e métodos tradicionais.

\subsection{Aprendizado por reforço}
\label{subsec:reforco}

No aprendizado por reforço, a alocação é formulada como um problema de decisão sequencial. Um agente observa o estado do sistema, como a utilização recente e o número de instâncias ativas, escolhe uma ação, como adicionar, remover ou manter instâncias, e recebe uma recompensa que penaliza o custo e as violações de SLA. Ao longo de muitas interações, o agente aprende uma política que maximiza a recompensa acumulada. Diferentemente da alocação guiada por previsão, previsão e decisão não são separadas: a política mapeia diretamente estados em ações, e a antecipação da demanda fica implícita no que o agente aprende. O aprendizado por reforço aparece tanto na classificação de \citeonline{lorido2014autoscaling} quanto na revisão de \citeonline{bodra2025ml}.

Essa revisão compara abordagens recentes de aprendizado de máquina para alocação de recursos em nuvem, incluindo métodos de aprendizado por reforço profundo, meta-heurísticas híbridas e modelos de previsão. Os autores observam que os trabalhos analisados são avaliados em ambientes, conjuntos de dados e métricas muito diferentes entre si, muitas vezes com cargas sintéticas, o que dificulta a comparação direta. Apontam ainda que os métodos de aprendizado por reforço profundo são sensíveis ao projeto da função de recompensa e aos hiperparâmetros e enfrentam limitações de escalabilidade conforme o espaço de estados cresce.

\textbf{Limitações e posicionamento.} Além dessa sensibilidade, uma política aprendida é difícil de interpretar: quando ela toma uma decisão ruim, não é simples identificar se o problema está na antecipação da demanda ou na regra de decisão. O treinamento também exige um ambiente de simulação e um grande número de episódios. Por esses motivos, este trabalho adota a abordagem desacoplada da Seção~\ref{subsec:alocacao-preditiva}, que permite isolar e medir o efeito da previsão sobre a alocação, e deixa o aprendizado por reforço fora do escopo (Seção~\ref{sec:escopo}). Comparar a abordagem desacoplada com uma política aprendida é uma extensão natural para trabalhos futuros.

\section{Comparativo dos trabalhos relacionados e lacuna}
\label{sec:lacuna}

A Tabela~\ref{tab:comparativo} resume os trabalhos mais próximos desta proposta segundo os critérios relevantes para o problema definido na Seção~\ref{sec:definicao-problema}: a técnica de previsão, os dados usados, as métricas de avaliação, se a previsão é acoplada a uma decisão de alocação e se diferentes previsores são comparados sob a mesma política de alocação.

\begin{table}[htbp]
\centering
\caption{Comparativo entre os trabalhos relacionados e a proposta deste trabalho.}
\label{tab:comparativo}
\footnotesize
\setlength{\tabcolsep}{3pt}
\begin{tabularx}{\textwidth}{
  >{\raggedright\arraybackslash}p{2.1cm}
  >{\raggedright\arraybackslash}X
  >{\raggedright\arraybackslash}p{1.9cm}
  >{\raggedright\arraybackslash}X
  >{\raggedright\arraybackslash}p{1.5cm}
  >{\raggedright\arraybackslash}p{1.7cm}}
\toprule
\textbf{Trabalho} &
\textbf{Técnica de previsão} &
\textbf{Dados} &
\textbf{Métricas} &
\textbf{Acopla à alocação} &
\textbf{Compara previsores sob a mesma política} \\
\midrule
\citeonline{calheiros2015arima} &
ARIMA &
Requisições à Wikipédia (1 h) &
Precisão, rejeições, violações de QoS e horas de máquina virtual &
Sim &
Não \\
\addlinespace
\citeonline{xu2022esdnn} &
GRU modificada com camadas convolucionais, multivariada &
Alibaba e Google &
MSE, RMSE, MAPE e máquinas ativas &
Parcial &
Não \\
\addlinespace
\citeonline{pan2023magicscaler} &
Atenção multiescala com processo gaussiano &
Alibaba Cloud (privados) &
MSE, MAE, custo e perda de QoS &
Sim &
Não \\
\addlinespace
\citeonline{rossi2023forecasting} &
Redes bayesianas (HBNN, LSTMD) &
Alibaba e Google &
Erro pontual e nível de serviço &
Não &
Não se aplica \\
\addlinespace
\citeonline{agullo2025enhancing} &
Tradicionais, aprendizado de máquina e profundo, com saída ajustada &
Alibaba e Google &
Erro, detecção de picos e superprovisionamento &
Não &
Não se aplica \\
\addlinespace
\citeonline{bodra2025ml} &
Revisão: aprendizado por reforço profundo, meta-heurísticas e previsores &
Variados, muitos sintéticos &
Heterogêneas entre os trabalhos &
Não se aplica &
Não se aplica \\
\addlinespace
\textbf{Este trabalho} &
GRU com saída direta (processador no PF~I; processador e memória no PF~II) e cinco métodos tradicionais &
Alibaba e Google &
Erro de previsão, custo, violações de SLA e ociosidade &
Sim &
Sim \\
\bottomrule
\end{tabularx}
\end{table}

A tabela evidencia duas linhas de trabalho que não se encontram. Os trabalhos de previsão \cite{xu2022esdnn, rossi2023forecasting, agullo2025enhancing} comparam vários modelos, mas por métricas calculadas sobre a própria previsão, em geral simétricas, que não distinguem a superestimação, que gera custo, da subestimação, que gera violações. \citeonline{rossi2023forecasting} mostram que essas métricas, isoladamente, não indicam a utilidade dos modelos para a gestão de recursos e apontam a integração entre previsão e alocação como questão em aberto, e \citeonline{agullo2025enhancing} mostram que previsores com bom erro médio falham justamente nos picos, quando as violações acontecem. Os trabalhos que acoplam a previsão à alocação \cite{calheiros2015arima, pan2023magicscaler} medem o efeito sobre o serviço, mas aplicam a decisão a um único previsor, comparado com políticas reativas, e o esDNN, que chega a simular o escalonamento, compara o resultado apenas com uma política baseada na média das observações anteriores \cite{xu2022esdnn}. A revisão de \citeonline{bodra2025ml} observa, por fim, que a literatura carece de comparações quantitativas sistemáticas, com ambientes, dados e métricas heterogêneos entre os trabalhos.

A lacuna que este trabalho busca preencher é, portanto, a de uma avaliação de ponta a ponta, com dados públicos de um \textit{cluster} real, que submeta previsores de complexidade crescente, do método ingênuo à rede neural, à mesma política de alocação, calibrada para o mesmo limite de violações $\varepsilon$, e meça o custo e as violações resultantes. A pergunta deixa de ser qual modelo erra menos e passa a ser quanto custo cada modelo economiza para um mesmo nível de serviço.

%======================================================================
\section{Tecnologias utilizadas}
\label{sec:tecnologias}

Esta seção descreve as linguagens, bibliotecas e ferramentas adotadas, organizadas conforme as etapas do \textit{pipeline}: tratamento dos dados, previsão de carga, simulação da alocação e ambiente de execução. O código do protótipo é escrito em Python 3.12, linguagem escolhida pela disponibilidade de bibliotecas maduras para séries temporais e aprendizado profundo, e está disponível em repositório público \cite{rosas2026repo}.

\subsection{Tratamento dos dados}
\label{subsec:tec-dados}

A tabela de uso de recursos por máquina do rastro da Alibaba de 2018 \cite{alibaba2018trace} ocupa cerca de 1,7~GB compactada. Para agregá-la em séries de demanda de processador e de memória com intervalos de cinco minutos sem carregar o arquivo inteiro na memória, é usado o DuckDB \cite{raasveldt2019duckdb}, um banco de dados analítico embutido que executa consultas na linguagem SQL (\textit{Structured Query Language}) diretamente sobre arquivos de valores separados por vírgula (CSV). As séries resultantes são gravadas em formato Parquet e manipuladas com pandas \cite{mckinney2010pandas} e NumPy \cite{harris2020numpy}. A mesma série alimenta todos os previsores, o que garante a comparabilidade entre eles.

\subsection{Modelos de previsão}
\label{subsec:tec-previsao}

Os métodos ingênuo, de média móvel e sazonal ingênuo são implementados diretamente com NumPy. O ARIMA e o Holt-Winters usam a biblioteca statsmodels \cite{seabold2010statsmodels} para a estimação dos parâmetros no conjunto de treino. A rede neural é implementada e treinada com PyTorch \cite{paszke2019pytorch}, que oferece flexibilidade para definir arquiteturas recorrentes e personalizar o ciclo de treinamento. Como a rede é pequena, é usada a versão do PyTorch que executa apenas na unidade central de processamento (CPU), o que dispensa uma unidade de processamento gráfico (GPU) para reproduzir os experimentos.

\subsection{Simulação}
\label{subsec:tec-simulacao}

O simulador adotado para a avaliação final é o CloudSim Plus \cite{silvafilho2017cloudsimplus}, um \textit{framework} em Java derivado do CloudSim \cite{calheiros2011cloudsim}. O CloudSim é um simulador de eventos discretos que modela centros de dados, servidores, máquinas virtuais e a execução de tarefas, incluindo políticas de compartilhamento dos núcleos de processamento entre máquinas virtuais e entre tarefas \cite{calheiros2011cloudsim}, e foi usado, por exemplo, para avaliar o provisionamento guiado por ARIMA de \citeonline{calheiros2015arima}. No Projeto Final II, ele receberá as previsões geradas em Python, como descrito na Seção~\ref{sec:solucao}.

Durante o Projeto Final I, o protótipo usou um simulador próprio, em Python, que implementa a antecedência da decisão $L$, as políticas e as métricas da Seção~\ref{subsec:metricas}; a troca de simulador é discutida na Seção~\ref{sec:mudancas-plano}. As figuras são geradas com Matplotlib \cite{hunter2007matplotlib}.

\subsection{Ambiente de execução e reprodutibilidade}
\label{subsec:tec-ambiente}

As dependências são gerenciadas com o uv, que fixa as versões de todas as bibliotecas em um arquivo de travamento, e o projeto pode ser executado também em um contêiner Docker. Um Makefile encadeia as etapas do \textit{pipeline}, de modo que todas as tabelas e figuras podem ser regeneradas a partir dos dados brutos com um único comando. Todos os parâmetros dos experimentos, como a divisão dos dados, os horizontes de previsão e as grades de margens, ficam em um arquivo de configuração, e nenhum é fixado no código. A correção dos componentes é verificada com testes automatizados escritos com pytest, e o código é versionado com Git no GitHub.